\documentclass[11pt]{article}
\usepackage{mathblatt}
\makeatletter
\setlength{\mbpfbe}{0mm}% keine Punkte (Bauregel 6.4)
\makeatother
% Rechenraum als Karo (Bauregel 6.9): n Rechenschritte = n mal 10 mm
\newcommand{\karo}[1]{\par\smallskip\noindent\begin{tikzpicture}
  \clip (0,0) rectangle ({\linewidth-1mm},{#1*10mm});
  \draw[black!16,line width=0.3pt,step=5mm] (0,0) grid ({\linewidth},{#1*10mm});
  \end{tikzpicture}\par}
\newcommand{\linie}[1]{\,{\color{black!45}\rule[-1pt]{#1}{0.4pt}}\,}
\newcommand{\pkt}{\textperiodcentered{}}
\begin{document}

\pfheftstil
\pflaufkopf[P10 \pkt{} mehr]{Seite oder Winkel mit sin, cos, tan \pkt{} Lösungen}{B8X}\pfseitevon
\pfheftkopf{Seite oder Winkel mit sin, cos, tan}{Lösungen \pkt{} Portion 2 von 2}
\begin{pfloesungb}{}
\lzb{1.}{}{AB $\approx$ 1118 m}{AB = 1400 · sin 53° $\approx$ 1118,1 m $\approx$ 1118 m}
\end{pfloesungb}
\begin{pfloesungb}{}
\lzb{2.}{}{AC $\approx$ 84,7 cm}{sin 24° = 34,47 : AC $\Rightarrow$ AC = 34,47 : sin 24° $\approx$ 84,7 cm}
\end{pfloesungb}
\begin{pfloesungb}{}
\lzb{3.}{}{tan $\alpha$ = r/s}{tan $\alpha$ = r/s}
\end{pfloesungb}
\begin{pfloesungb}{}
\lzb{4.}{}{$\alpha$ $\approx$ 65°}{cos $\alpha$ = 5 : 12 $\Rightarrow$ $\alpha$ = cos$^{-1}$($\tfrac{5}{12}$) $\approx$ 65,4° $\approx$ 65°}
\end{pfloesungb}
\begin{pfloesungb}{}
\lzb{5.}{}{sin $\alpha$ = r/t}{sin $\alpha$ = r/t}
\end{pfloesungb}
\begin{pfloesungb}{}
\lzb{6.}{}{sin $\beta$ = r/t}{sin $\beta$ = r/t}
\end{pfloesungb}
\begin{pfloesungb}{}
\lzb{7.}{}{PB $\approx$ 1,85 m}{PB = 1,35 : tan 36,1° $\approx$ 1,85 m}
\end{pfloesungb}
\end{document}